\documentclass[uplatex,10pt,twocolumn,a4paper]{article}

% -------------------------------------------------------
% Packages
% -------------------------------------------------------
\usepackage[dvipdfmx]{graphicx}
\usepackage[dvipdfmx,hidelinks]{hyperref}
\usepackage{amsmath,amssymb,amsthm}
\usepackage{booktabs}
\usepackage{geometry}
\usepackage{algorithm}
\usepackage{algorithmic}
\usepackage{xcolor}
\usepackage{caption}
\usepackage{subcaption}
\usepackage{cite}
\usepackage{url}

\geometry{a4paper, left=18mm, right=18mm, top=22mm, bottom=22mm,
          columnsep=6mm}

\newcommand{\TODO}[1]{\textcolor{red}{\textbf{[#1]}}}
\newcommand{\ie}{\textit{i.e.,}\ }
\newcommand{\eg}{\textit{e.g.,}\ }
\newcommand{\etal}{\textit{et al.}}

% -------------------------------------------------------
% Title
% -------------------------------------------------------
\title{
  \vspace{-6mm}
  \Large\textbf{Context-Aware Preference Tuning for Masked Scene Text Prediction\\
  with Agentic Self-Correction}
}

\author{
  Sho Yamamoto\\
  \small Hitotsubashi University\\
  \small \texttt{5123062h@g.hit-u.ac.jp}
}

\date{}

% -------------------------------------------------------
\begin{document}
% -------------------------------------------------------

\maketitle
\thispagestyle{empty}

% -------------------------------------------------------
\begin{abstract}
% -------------------------------------------------------
We address the task of \emph{masked scene text prediction}: given a
natural image in which a text region has been occluded together with a
binary mask indicating the occluded area, a model must recover the
original text string from visual context alone.
Existing Vision-Language Models (VLMs) can be adapted to this task via
supervised fine-tuning, but standard cross-entropy training treats all
plausible outputs equally and fails to exploit the rich contextual
information available from surrounding scene text.

We propose a two-phase fine-tuning pipeline for Qwen3-VL-2B-Instruct.
\textbf{Phase~1} performs supervised fine-tuning (SFT) with paraphrase
augmentation, teaching the model to predict masked text from
(masked image, binary mask) pairs.
\textbf{Phase~2} applies \emph{context-aware online Direct Preference
Optimisation} (DPO): at each step, multiple candidate predictions are
sampled from the current policy and ranked by their embedding-space
similarity to the surrounding scene context---both the neighbouring
paragraph texts and the masked image---using a visual-language embedding
model.
This automatically constructs preference pairs without any human
annotation, directly encouraging predictions that are semantically
consistent with the scene.
We further embed the fine-tuned model into a \textbf{self-correcting
agent loop}: the model evaluates rendered text across two visual
granularities and iteratively adapts generation parameters to improve
both character accuracy and background harmony.

On the HierText validation set, our two-phase fine-tuned model achieves
\TODO{XX}\% CER reduction and \TODO{XX}-point ANLS improvement over the
base Qwen3-VL-2B model.
Ablation studies confirm that context-aware DPO contributes
\TODO{XX}~CER points beyond SFT alone.

\end{abstract}

% -------------------------------------------------------
\section{Introduction}
\label{sec:intro}
% -------------------------------------------------------

Reading text that has been partially or fully hidden in a scene image
is a problem that arises naturally in document restoration, copyright
enforcement, and augmented-reality overlay verification.
Unlike standard scene text recognition (STR), where the text is
\emph{visible}, masked text prediction requires the model to infer the
hidden content from surrounding visual and textual context: the shape
of the occluded region, the style of visible neighbouring text, and the
semantic coherence of the scene.

Recent Vision-Language Models such as the Qwen-VL
family~\cite{qwen2vl2024} possess broad visual understanding that makes
them natural candidates for this task.
However, standard supervised fine-tuning (SFT) via cross-entropy loss
has two shortcomings.
First, it provides no mechanism to prefer outputs that are
\emph{contextually plausible} over those that are merely character-wise
plausible.
Second, when multiple reasonable phrasings exist (\eg through
capitalisation, punctuation, or synonym variation), SFT conflates them
by penalising all non-gold-standard predictions equally.

We address both limitations by combining SFT with \emph{context-aware
online Direct Preference Optimisation}~\cite{dpo2023}.
Our key insight is that the correct hidden text should be semantically
consistent with the surrounding scene: nearby paragraph texts share
topic and style, and the masked image encodes spatial context.
Concretely, at each training step we sample $N$ candidate predictions,
score them by their cosine similarity to a context embedding formed from
surrounding paragraph texts and the masked image via a VL embedding
model, and use the best/worst pair as the chosen/rejected response for
a DPO update.
This requires \emph{no human preference labels} and can run online
without a separate reward model.

Beyond the prediction model, we embed it in a \textbf{self-correcting
agent loop} for scene text rendering.
Given a target text, the agent iteratively generates candidate
inpaintings with a diffusion-based text renderer, evaluates rendered
quality at two visual granularities using the fine-tuned VLM, and feeds
structured parameter corrections back to the renderer.
A \emph{scene-anchoring} mechanism stabilises the background caption
across iterations, preventing prompt drift.

Our contributions are:

\begin{itemize}
  \item \textbf{Two-phase SFT + context-aware online DPO.}
    We show that DPO with embedding-based, annotation-free preference
    construction improves masked text prediction beyond SFT alone.

  \item \textbf{Paraphrase augmentation.}
    We pre-generate lexical paraphrases of ground-truth labels and
    inject them randomly during SFT, improving robustness to surface
    variation.

  \item \textbf{Context-aware preference scoring.}
    Candidate responses are ranked by cosine similarity to a joint
    text-and-image context embedding, directly leveraging scene
    coherence as a supervision signal.

  \item \textbf{Self-correcting rendering agent.}
    The fine-tuned VLM evaluates rendered text at two spatial
    granularities and provides actionable parameter corrections in a
    structured agent loop.
\end{itemize}

% -------------------------------------------------------
\section{Related Work}
\label{sec:related}
% -------------------------------------------------------

\paragraph{Scene text recognition.}
Classic pipelines combine CNNs with CTC or attention-based
decoders~\cite{crnn2015, aster2019, abinet2021}.
TrOCR~\cite{trocr2021} adapts encoder-decoder transformers pre-trained
on image-text pairs, achieving strong STR performance.
LLaVAR~\cite{llavar2023} shows that VLMs can be instruction-tuned for
text-rich image understanding.
None of these works addresses \emph{masked} text, where character
appearance is unavailable and context reasoning is essential.

\paragraph{Direct Preference Optimisation.}
DPO~\cite{dpo2023} reframes RLHF as a classification objective over
chosen/rejected response pairs, eliminating the need for an explicit
reward model.
Online DPO~\cite{onlinedpo2024} generates preference pairs from the
current policy, reducing distribution mismatch.
We extend online DPO to the multimodal masked-text setting and replace
human annotations with embedding-based scene-coherence scoring.

\paragraph{Preference learning without human labels.}
RLAIF~\cite{rlaif2023} uses an LLM as the preference annotator.
Self-Play Fine-Tuning (SPIN)~\cite{spin2024} constructs pairs by
treating the SFT model's own outputs as rejected responses.
Our approach is orthogonal: we use a separate embedding model to score
candidate predictions by scene-context similarity, which is a
domain-specific signal unavailable to general-purpose LLM annotators.

\paragraph{Text-in-image generation.}
AnyText~\cite{anytext2023} conditions a latent diffusion model on text
glyphs and auxiliary guidance for text inpainting.
TextDiffuser~\cite{textdiffuser2023} and
GlyphControl~\cite{glyphcontrol2023} add layout control.
These models operate in a single forward pass with no quality feedback;
our agent loop closes this gap.

\paragraph{Agentic self-refinement.}
ReAct~\cite{react2022} and Reflexion~\cite{reflexion2023} equip LLMs
with iterative self-correction.
Self-Refine~\cite{selfrefine2023} applies iterative feedback to text
generation.
We extend this paradigm to multimodal scene text rendering, with a
VLM providing structured visual feedback at two granularities.

% -------------------------------------------------------
\section{Method}
\label{sec:method}
% -------------------------------------------------------

\subsection{Problem Formulation}
\label{sec:problem}

Let $I \in \mathbb{R}^{H \times W \times 3}$ be a natural scene image
and $M \in \{0,1\}^{H \times W}$ a binary mask indicating a text
region ($M_{ij}=1$ at text pixels).
The masked image $\tilde{I} = I \odot (1 - M)$ has the text region
set to white.
We also observe a set of surrounding paragraph texts
$\mathcal{C} = \{c_1, \ldots, c_K\}$ from the same image.

The \textbf{masked text prediction} task is: given $(\tilde{I}, M)$,
predict the ground-truth text string $y^* = (y_1,\ldots,y_T)$ that
occupied the masked region.
The \textbf{scene text rendering} task (Section~\ref{sec:stage2}) is:
given $(\tilde{I}, M, y^*)$, produce a composited image $\hat{I}$ that
renders $y^*$ naturally into the scene.

\subsection{Stage 1: Supervised Fine-Tuning}
\label{sec:sft}

\paragraph{Model and adaptation.}
We fine-tune Qwen3-VL-2B-Instruct~\cite{qwen3vl2025} using
QLoRA~\cite{qlora2023} with 4-bit NF4 quantisation and double
quantisation~\cite{qlora2023}.
LoRA adapters (rank $r=4$, scaling $\alpha=16$, dropout $p=0.05$) are
attached to the query and value projection matrices ($\mathbf{W}_q$,
$\mathbf{W}_v$) of every attention layer.

\paragraph{Input format.}
Each training sample presents two images to the VLM in sequence:
(1) the masked scene image $\tilde{I}$, and
(2) the binary mask $M$ converted to a three-channel grayscale image.
The system prompt is fixed across all samples:
\begin{quote}
  \small
  \textit{``You are given two images: (1) a masked image where text
  regions are hidden; (2) a binary mask where the white region is text
  and the black region is background.
  Task: Predict the text in the white masked region.''}
\end{quote}
The assistant response is the ground-truth text $y^*$.
We apply the Qwen3 chat template with \texttt{add\_generation\_prompt=False}
so that both the prompt and the response are encoded in a single forward
pass, and set labels to $-100$ for all prompt tokens so that the loss
is computed only over the response.

\paragraph{Paraphrase augmentation.}
A fixed ground-truth label during SFT risks overfitting to a single
surface form.
We pre-generate lexical paraphrases for each training sample
(capitalisation variants, punctuation variants, minor rephrasing) and
cache them.
With probability $0.5$ per sample, the target text is replaced by a
uniformly sampled paraphrase.
This acts as a form of label smoothing tailored to the character-level
metric space and reduces sensitivity to case or punctuation ambiguity.

\paragraph{Data augmentation.}
Both $\tilde{I}$ and $M$ undergo: random rotation by $\theta \sim
\mathcal{U}(-5^\circ, 5^\circ)$ (applied with the same angle to both
images to preserve alignment); colour jitter (brightness $\pm 0.3$,
contrast $\pm 0.3$, saturation $\pm 0.2$, hue $\pm 0.05$) applied to
$\tilde{I}$ only; and Gaussian blur (radius $r \sim \mathcal{U}(0.5,
1.5)$) applied to $\tilde{I}$ with probability $0.5$.

\paragraph{Sample selection.}
HierText provides hierarchical word/line/paragraph annotations.
We use only paragraphs with at most 9 words to keep the prediction
tractable.
For each image we retain at most 3 paragraphs, ranked by a
coverage-weighted score:
\begin{equation}
  \text{score}(p) = \rho(p) \cdot \log\bigl(1 + n_w(p)\bigr),
  \label{eq:score}
\end{equation}
where $\rho(p)$ is the convex-hull area of paragraph $p$ as a fraction
of image area and $n_w(p)$ is the word count.
This prioritises spatially prominent paragraphs with sufficient text.

\paragraph{Training objective.}
The SFT loss is the standard causal language modelling cross-entropy:
\begin{equation}
  \mathcal{L}_\mathrm{SFT} = -\sum_{t=1}^{T} \log P_\theta(y_t \mid
  \tilde{I}, M, y_{<t}),
  \label{eq:sft}
\end{equation}
computed only over response tokens.

\subsection{Stage 2: Context-Aware Online DPO}
\label{sec:dpo}

After SFT the model is further tuned with online DPO to incorporate
scene-context preferences.
At each optimisation step, a subset of batch samples undergoes the
preference construction procedure described below.

\paragraph{Candidate generation.}
For each selected sample $i$, we sample $N=3$ candidate responses
$\{\hat{y}^{(1)}, \ldots, \hat{y}^{(N)}\}$ from the current policy
$\pi_\theta$ in inference mode, using temperature $T=0.9$ and a maximum
of 32 new tokens.
Duplicate candidates are de-duplicated; if fewer than 2 unique
candidates remain, the sample is skipped.

\paragraph{Context embedding.}
We use a visual-language embedding model (Qwen3-VL-Embedding) to
construct a context vector $\mathbf{v}_\mathcal{C}$ for sample $i$.
Surrounding paragraph texts $\mathcal{C} = \{c_1, \ldots, c_K\}$
(up to $K=8$) are each encoded and averaged:
\begin{equation}
  \mathbf{v}_\text{text} = \ell_2\text{-norm}\!\left(
    \frac{1}{|\mathcal{C}|}\sum_{k} f_\phi(c_k)
  \right) \in \mathbb{R}^d,
  \label{eq:text_embed}
\end{equation}
where $f_\phi$ denotes the embedding model's last-token hidden state.
Additionally, if the masked image path $\tilde{I}_i$ is available, an
image embedding is extracted:
\begin{equation}
  \mathbf{v}_\text{img} = \ell_2\text{-norm}\!\left(
    \frac{1}{L}\sum_{\ell=1}^{L} f_\phi(\tilde{I}_i)_\ell
  \right) \in \mathbb{R}^d,
  \label{eq:img_embed}
\end{equation}
where the mean is over sequence positions $\ell$.
The joint context embedding is
\begin{equation}
  \mathbf{v}_\mathcal{C} = \ell_2\text{-norm}
    (\mathbf{v}_\text{text} + \mathbf{v}_\text{img}).
  \label{eq:joint_embed}
\end{equation}
When no image is available, $\mathbf{v}_\mathcal{C} = \mathbf{v}_\text{text}$.

\paragraph{Preference construction.}
Each candidate $\hat{y}^{(n)}$ is encoded as a text embedding
$\mathbf{u}^{(n)} = \ell_2\text{-norm}(f_\phi(\hat{y}^{(n)}))$ and
scored by cosine similarity to the context:
\begin{equation}
  s^{(n)} = \mathbf{u}^{(n)\top} \mathbf{v}_\mathcal{C}.
  \label{eq:score_dpo}
\end{equation}
The candidate with the highest score is designated
\emph{chosen} ($y_+$) and the one with the lowest score is
\emph{rejected} ($y_-$).

\paragraph{DPO loss.}
We use the standard DPO objective~\cite{dpo2023} with the base model
(LoRA adapters disabled) as the reference policy $\pi_\mathrm{ref}$:
\begin{equation}
  \mathcal{L}_\mathrm{DPO} = -\log\sigma\!\Bigl(
    \beta \bigl[
      \log\tfrac{\pi_\theta(y_+ \mid \tilde{I}, M)}
                {\pi_\mathrm{ref}(y_+ \mid \tilde{I}, M)}
      -
      \log\tfrac{\pi_\theta(y_- \mid \tilde{I}, M)}
                {\pi_\mathrm{ref}(y_- \mid \tilde{I}, M)}
    \bigr]
  \Bigr),
  \label{eq:dpo}
\end{equation}
where $\beta=0.1$ controls the deviation from the reference policy.
Sequence log-probabilities are computed as the sum of per-token
log-softmax values over the response span.

\paragraph{Combined training loss.}
During the DPO phase the cross-entropy loss from Eq.~\ref{eq:sft}
is retained and the DPO term is added with weight $w=0.1$:
\begin{equation}
  \mathcal{L} = \mathcal{L}_\mathrm{SFT} +
                w \cdot \mathcal{L}_\mathrm{DPO}.
  \label{eq:total}
\end{equation}
To avoid excessive overhead, the DPO update is applied every
$\Delta_\mathrm{DPO} = 100$ gradient steps, and 2 samples per batch
are selected for preference construction at each application.

\subsection{Stage 3: Self-Correcting Rendering Agent}
\label{sec:stage2}

The fine-tuned VLM is embedded in a self-correcting agent loop for
scene text rendering (Algorithm~\ref{alg:agent}).
The loop uses AnyText~\cite{anytext2023} as the text inpainting
backbone and the VLM as a dual-granularity evaluator.

\begin{algorithm}[t]
  \caption{TextCorrectionAgent.run($\tilde{I}$, $M$, $y^*$)}
  \label{alg:agent}
  \begin{algorithmic}[1]
    \STATE $y \leftarrow y^*$;\ $\text{scene} \leftarrow \emptyset$;\
           $\text{style} \leftarrow \emptyset$;\ $i \leftarrow 0$
    \WHILE{$i < K$}
      \STATE $i \leftarrow i + 1$
      \STATE $\hat{I} \leftarrow \textsc{AnyText}(y,\, M,\, \tilde{I},\,
             [\text{scene};\, \text{style}])$
      \STATE $e \leftarrow \textsc{Evaluate}(\hat{I},\, y^*,\, M)$
             \COMMENT{Sec.~\ref{sec:evaluator}}
      \IF{$e.s \geq \tau$ \textbf{or} \textsc{NoImprovement}($3$)}
        \RETURN $\arg\max_j \hat{I}_j$
      \ENDIF
      \STATE update rendering params from $e$
      \IF{$\text{scene} = \emptyset$}
        \STATE $\text{scene} \leftarrow e.\texttt{suggested\_prompt}$
      \ENDIF
      \STATE $\text{style} \leftarrow e.\texttt{text\_style\_instruction}$
    \ENDWHILE
    \RETURN $\arg\max_j \hat{I}_j$
  \end{algorithmic}
\end{algorithm}

\paragraph{Dual-granularity evaluator.}
\label{sec:evaluator}
Given a generated image $\hat{I}$, the VLM receives two inputs:

\begin{enumerate}
  \item \textbf{Text-region crop.}
    The bounding box of mask $M$ is expanded by margin
    $m = \max(\max(w, h),\, 50)$ pixels on each side,
    square-cropped, and upscaled to $768\times768$.
    This enables character-level legibility assessment.
  \item \textbf{Full composite image.}
    The complete $\hat{I}$ enables holistic background-harmony scoring
    and scene captioning.
\end{enumerate}

The VLM returns a structured response containing:
text accuracy $s_\mathrm{text} \in [0,1]$,
background harmony $s_\mathrm{harmony} \in [0,1]$,
overall score $s = (s_\mathrm{text} + s_\mathrm{harmony})/2$,
a background scene caption (\texttt{suggested\_prompt}),
a rendering style instruction (\texttt{text\_style\_instruction}),
and optional parameter adjustments
$\Delta s_\mathrm{cfg}$, $\Delta r_\mathrm{mask}$,
$\Delta r_\mathrm{crop}$.

\paragraph{Structured parameter adaptation.}
Three AnyText parameters are updated each iteration:
$s_\mathrm{cfg}$ (classifier-free guidance scale),
$r_\mathrm{mask}$ (mask dilation factor), and
$r_\mathrm{crop}$ (crop margin).
Table~\ref{tab:rules} lists the failure-mode-to-correction mapping
output by the evaluator.

\begin{table}[t]
  \centering
  \caption{Parameter adaptation rules output by the evaluator.}
  \label{tab:rules}
  \small
  \begin{tabular}{lc}
    \toprule
    Observed failure & Corrective action \\
    \midrule
    Characters cramped / overlapping & $r_\mathrm{mask} > 1.0$ \\
    Characters too sparse            & $r_\mathrm{mask} < 1.0$ \\
    Excess background in crop        & $r_\mathrm{crop} < 1.0$ \\
    Text too long for region         & shorten $y$ \\
    Garbled characters               & $s_\mathrm{cfg} \in [13,15]$ \\
    Artefacts / white blobs          & $s_\mathrm{cfg} \in [7.5,\,9]$ \\
    \bottomrule
  \end{tabular}
\end{table}

\paragraph{Scene anchoring.}
The background caption produced on the first iteration is \emph{fixed}
for all subsequent iterations; only the style instruction is updated.
This prevents the evaluator from inadvertently describing rendered
characters as part of the scene, which would pollute subsequent
generation prompts.
The generation prompt at iteration $i$ is:
\begin{equation}
  p_i = [\text{scene};\, \text{style}_i],
  \label{eq:prompt}
\end{equation}
where $[\cdot;\cdot]$ denotes comma-separated concatenation.

\paragraph{Early stopping.}
The loop terminates when $s \geq \tau = 0.9$ (success) or when the
overall score does not improve over $N_\mathrm{patience} = 3$
consecutive iterations.
The maximum iteration count is $K = 5$.

% -------------------------------------------------------
\section{Experiments}
\label{sec:experiments}
% -------------------------------------------------------

\subsection{Setup}

\paragraph{Dataset.}
We use HierText~\cite{hiertext2022}, a large-scale dataset with
hierarchical word/line/paragraph annotations on natural scene images.
After filtering to paragraphs with $\leq 9$ words and applying the
coverage-weighted selection (Eq.~\ref{eq:score}), we obtain
\TODO{N\_train} training and \TODO{N\_val} validation samples.

\paragraph{Evaluation metrics.}
\begin{itemize}
  \item \textbf{CER} (Character Error Rate): $\text{edit}(\hat{y},
    y^*) / |y^*|$. Lower is better.
  \item \textbf{NED} (Normalised Edit Distance): $\text{edit}(\hat{y},
    y^*) / \max(|\hat{y}|, |y^*|)$. Lower is better.
  \item \textbf{ANLS}: $1 - \text{NED}$. Higher is better.
  \item \textbf{Acc@1}: Exact-match accuracy after stripping
    whitespace. Higher is better.
\end{itemize}

\paragraph{Training details.}
SFT runs for 6 epochs with AdamW (paged 8-bit), learning rate
$\eta = 2 \times 10^{-5}$, cosine schedule, $10\%$ linear warmup, and
per-device batch size 4.
The DPO phase uses the same optimiser and schedule; the DPO-specific
hyperparameters are $\beta = 0.1$, $w = 0.1$, $N = 3$ candidates,
$\Delta_\mathrm{DPO} = 100$ steps.
All experiments use a single GPU with 4-bit NF4 QLoRA.

\subsection{Masked Text Prediction}

Table~\ref{tab:ocr} compares masked text prediction accuracy.

\begin{table}[t]
  \centering
  \caption{Masked text prediction on HierText validation.}
  \label{tab:ocr}
  \small
  \begin{tabular}{lcccc}
    \toprule
    Model & CER$\downarrow$ & NED$\downarrow$ & ANLS$\uparrow$ & Acc@1$\uparrow$ \\
    \midrule
    Qwen3-VL-2B (base)         & \TODO{.XX} & \TODO{.XX} & \TODO{.XX} & \TODO{.XX} \\
    + SFT                      & \TODO{.XX} & \TODO{.XX} & \TODO{.XX} & \TODO{.XX} \\
    + SFT + paraphrase aug.    & \TODO{.XX} & \TODO{.XX} & \TODO{.XX} & \TODO{.XX} \\
    + SFT + context-aware DPO  & \TODO{.XX} & \TODO{.XX} & \TODO{.XX} & \TODO{.XX} \\
    \textbf{Ours (full)}       & \TODO{.XX} & \TODO{.XX} & \TODO{.XX} & \TODO{.XX} \\
    \bottomrule
  \end{tabular}
\end{table}

\paragraph{Effect of context-aware DPO.}
Context-aware DPO improves over SFT alone by encouraging predictions
that are semantically coherent with the surrounding scene.
The benefit is most pronounced for ambiguous samples where multiple
character sequences are visually plausible but differ in scene-level
coherence (\eg product names matching visible brand text, or numbers
consistent with surrounding signage).

\paragraph{Effect of paraphrase augmentation.}
Paraphrase augmentation provides gains on samples with alternative
valid surface forms, reducing CER on mixed-case and punctuated text.
The combination of augmentation and DPO is complementary: augmentation
addresses surface-form robustness during SFT, while DPO addresses
semantic consistency during preference tuning.

\subsection{Ablation: DPO Scoring Strategy}

Table~\ref{tab:ablation_dpo} compares context-aware scoring against
alternative preference construction strategies.

\begin{table}[t]
  \centering
  \caption{Ablation on DPO preference construction.}
  \label{tab:ablation_dpo}
  \small
  \begin{tabular}{lcc}
    \toprule
    Scoring strategy & CER$\downarrow$ & ANLS$\uparrow$ \\
    \midrule
    SFT only (no DPO)            & \TODO{.XX} & \TODO{.XX} \\
    DPO w/ text context only     & \TODO{.XX} & \TODO{.XX} \\
    DPO w/ image context only    & \TODO{.XX} & \TODO{.XX} \\
    DPO w/ text + image (ours)   & \TODO{.XX} & \TODO{.XX} \\
    \bottomrule
  \end{tabular}
\end{table}

Combining textual and visual context yields the best performance,
confirming that surrounding paragraph texts and the masked image
encode complementary information for scoring candidate predictions.

\subsection{Agent Loop: Scene Text Rendering Quality}

Table~\ref{tab:rendering} reports VLM-Score and convergence rate for
the self-correcting agent against baselines.

\begin{table}[t]
  \centering
  \caption{Scene text rendering quality on HierText validation.}
  \label{tab:rendering}
  \small
  \begin{tabular}{lcc}
    \toprule
    Method & VLM-Score$\uparrow$ & Conv.\ rate$\uparrow$ \\
    \midrule
    AnyText (single shot)           & \TODO{0.XX} & --- \\
    AnyText + OCR feedback          & \TODO{0.XX} & \TODO{XX}\% \\
    Ours (agent, $K=5$, base VLM)  & \TODO{0.XX} & \TODO{XX}\% \\
    Ours (agent, $K=5$, tuned VLM) & \TODO{0.XX} & \TODO{XX}\% \\
    \bottomrule
  \end{tabular}
\end{table}

Using the fine-tuned VLM as the evaluator improves over the base model,
demonstrating that HierText-specific knowledge acquired during SFT+DPO
transfers to the evaluation role.
Scene anchoring prevents caption drift and contributes
\TODO{X} VLM-Score points in the ablation (Table~\ref{tab:ablation_agent}).

\begin{table}[t]
  \centering
  \caption{Ablation on agent components (VLM-Score).}
  \label{tab:ablation_agent}
  \small
  \begin{tabular}{lcc}
    \toprule
    Configuration & VLM-Score & Conv.\ rate \\
    \midrule
    Full agent                       & \TODO{0.XX} & \TODO{XX}\% \\
    w/o scene anchoring              & \TODO{0.XX} & \TODO{XX}\% \\
    w/o dual-granularity evaluation  & \TODO{0.XX} & \TODO{XX}\% \\
    w/o parameter adaptation         & \TODO{0.XX} & \TODO{XX}\% \\
    w/o style instruction            & \TODO{0.XX} & \TODO{XX}\% \\
    \bottomrule
  \end{tabular}
\end{table}

\begin{figure}[t]
  \centering
  \fbox{\rule{0pt}{30mm}\rule{0.96\linewidth}{0pt}}
  \caption{
    Mean VLM-Score across agent iterations (shaded: $\pm$1~std).
    The dashed line marks the success threshold $\tau = 0.9$.
  }
  \label{fig:convergence}
\end{figure}

% -------------------------------------------------------
\section{Discussion}
\label{sec:discussion}
% -------------------------------------------------------

\paragraph{Annotation-free preference learning.}
Our context-aware DPO requires only the existing HierText paragraph
annotations for the surrounding-text context.
No additional human preference labels are needed.
This makes the approach directly applicable to any dataset with
hierarchical or spatially clustered text annotations.

\paragraph{Model size and efficiency.}
We deliberately chose the 2B parameter model over larger alternatives
to demonstrate that strong masked-text prediction is achievable with a
compact VLM under QLoRA.
The full SFT+DPO training fits on a single GPU with 4-bit
quantisation.
The online DPO candidate generation (3 sequences per selected sample,
every 100 steps) adds approximately \TODO{XX}\% wall-clock overhead
relative to SFT alone.

\paragraph{Limitations.}
The embedding-based preference scoring relies on the assumption that
the correct text is more semantically similar to surrounding context
than incorrect alternatives.
This assumption may fail for isolated signs, foreign-language text, or
numerals with no topical connection to surrounding paragraphs.
The agent loop additionally incurs $O(K)$ generator calls per image,
which limits throughput.
Extension to other inpainting backbones beyond AnyText is
straightforward given the modular architecture.

% -------------------------------------------------------
\section{Conclusion}
\label{sec:conclusion}
% -------------------------------------------------------

We presented a two-phase fine-tuning approach for masked scene text
prediction.
Phase~1 (SFT with paraphrase augmentation) establishes a strong
baseline by training Qwen3-VL-2B-Instruct on (masked image, binary
mask) pairs.
Phase~2 (context-aware online DPO) further improves predictions by
scoring candidates using embedding-space similarity to the surrounding
scene context---both paragraph texts and the masked image---without any
human preference labels.
The fine-tuned model is embedded in a self-correcting agent loop that
iteratively adjusts AnyText rendering parameters guided by structured
dual-granularity visual evaluation.
On HierText, the full system achieves \TODO{XX}\% CER reduction over
the base model and \TODO{XX} VLM-Score points improvement over
single-shot rendering.
We believe annotation-free preference construction via scene context is
a broadly applicable technique for other structured text-in-image tasks.

% -------------------------------------------------------
\bibliographystyle{plain}
\begin{thebibliography}{99}

\bibitem{qwen3vl2025}
Qwen Team,
``Qwen3-VL: Seeing and Reading the World,''
\textit{arXiv:2504.10479}, 2025.

\bibitem{qwen2vl2024}
Qwen Team,
``Qwen2-VL: Enhancing Vision-Language Model's Perception of the World
at Any Resolution,''
\textit{arXiv:2409.12191}, 2024.

\bibitem{hiertext2022}
A.~Long \etal,
``Towards End-to-End Unified Scene Text Detection and Layout Analysis,''
\textit{CVPR}, 2022.

\bibitem{dpo2023}
R.~Rafailov \etal,
``Direct Preference Optimization: Your Language Model is Secretly a
Reward Model,''
\textit{NeurIPS}, 2023.

\bibitem{onlinedpo2024}
T.~Guo \etal,
``Direct Language Model Alignment from Online AI Feedback,''
\textit{arXiv:2402.04792}, 2024.

\bibitem{qlora2023}
T.~Dettmers \etal,
``QLoRA: Efficient Finetuning of Quantized LLMs,''
\textit{NeurIPS}, 2023.

\bibitem{anytext2023}
T.~Tuo \etal,
``AnyText: Multilingual Visual Text Generation and Editing,''
\textit{arXiv:2311.03054}, 2023.

\bibitem{textdiffuser2023}
J.~Chen \etal,
``TextDiffuser: Diffusion Models as Text Painters,''
\textit{NeurIPS}, 2023.

\bibitem{glyphcontrol2023}
Y.~Yang \etal,
``GlyphControl: Glyph Conditional Control for Visual Text Generation,''
\textit{NeurIPS}, 2023.

\bibitem{crnn2015}
B.~Shi \etal,
``An End-to-End Trainable Neural Network for Image-based Sequence
Recognition,''
\textit{TPAMI}, 2017.

\bibitem{aster2019}
B.~Shi \etal,
``ASTER: An Attentional Scene Text Recognizer with Flexible
Rectification,''
\textit{TPAMI}, 2019.

\bibitem{abinet2021}
F.~Fang \etal,
``Read Like Humans: Autonomous, Bidirectional and Iterative Language
Modeling for Scene Text Recognition,''
\textit{CVPR}, 2021.

\bibitem{trocr2021}
M.~Li \etal,
``TrOCR: Transformer-based Optical Character Recognition with
Pre-trained Models,''
\textit{AAAI}, 2023.

\bibitem{llavar2023}
Y.~Zhang \etal,
``LLaVAR: Enhanced Visual Instruction Tuning for Text-Rich Image
Understanding,''
\textit{arXiv:2306.17107}, 2023.

\bibitem{rlaif2023}
H.~Lee \etal,
``RLAIF: Scaling Reinforcement Learning from Human Feedback with AI
Feedback,''
\textit{arXiv:2309.00267}, 2023.

\bibitem{spin2024}
Z.~Chen \etal,
``Self-Play Fine-Tuning Converts Weak Language Models to Strong Language
Models,''
\textit{ICML}, 2024.

\bibitem{react2022}
S.~Yao \etal,
``ReAct: Synergizing Reasoning and Acting in Language Models,''
\textit{ICLR}, 2023.

\bibitem{reflexion2023}
N.~Shinn \etal,
``Reflexion: Language Agents with Verbal Reinforcement Learning,''
\textit{NeurIPS}, 2023.

\bibitem{selfrefine2023}
A.~Madaan \etal,
``Self-Refine: Iterative Refinement with Self-Feedback,''
\textit{NeurIPS}, 2023.

\end{thebibliography}

\end{document}
